\documentstyle[%


righttag%


,11pt%


,amssymb%


,russian%               remove in Eng.


,izvru%                 Set izven% in Eng.


]{amsart}





% 





\newcommand{\rank}{\operatorname{rank}}


\newcommand{\tts}[1]{}





\theoremstyle{definition}


\theorembodyfont{\rm}


\newtheorem{cornonum}{\hskip\parindent Следствие}


\renewcommand{\thecornonum}{}


\renewcommand{\baselinestretch}{0.99}





%\hoffset=-15mm%                  For Eng.


\setcounter{page}{3}


\god{2004}


\nomer{10~(509)} %                        Remove in Eng.


%\nomer{Vol.\,48, No.\,10}%               Set in Eng.


%\str{\pageref{begin}--\pageref{end}}%  Set in Eng.


\udk{514.752}





\author{\it М.А.\,Акивис, В.В.\,Гольдберг, А.В.\,Чакмазян}


% NB:   ^^ remove in Eng.


\title{Индуцированные связности на многообразиях\\


в пространствах с фундаментальными группами}





\date{}


%\pagestyle{myheadings}%         Set in Eng.


\pagestyle{plain} %              Remove in Eng.


\begin{document}


%\thispagestyle{plain}%          Set in Eng.


\maketitle


\label{begin}









\section*{Введение}



Теория конгруэнций и псевдоконгруэнций подпространств проективного


пространства тесно связана с теорией многообразий с вырожденным


гауссовым отображением.





Теория конгруэнций в трехмерном проективном


пространстве $\Bbb{P}^3$, а также в


трехмерных пространствах, снабженных проективной структурой (таких как


аффинное, евклидово и неевклидовы пространства), изучалась многими


геометрами.


Имеются обширные монографии, посвященные этому предмету


(напр., \cite{Fi50}).



В данной работе установлены взаимосвязи теории многообразий


в проективных пространствах с вырожденным гауссовым отображением


с теориeй конгруэнций и псевдоконгруэнций подпространств и показано,


как эти две теории могут быть применены к построению индуцированных


связностей на подмногообразиях в проективных пространствах и других


пространствах, снабженных проективной структурой.








\section{Основные уравнения многообразия \protect\\ с вырожденным гауссовым


отображением}





Многообразие $X$ размерности $n$, которое является подмногообразием


проективного пространства $\Bbb{P}^N$, $X\subset \Bbb{P}^N$,


называется {\em тангенциально вырожденным


многообразием} или {\em многообразием с вырожденным гауссовым


отображением}, если ранг его гауссова отображения


$$


\gamma: X \rightarrow \Bbb{G} (n, N)


$$
меньше, чем $n$, $ 0 \leq r = \rank \gamma < n$; здесь


$\gamma (x) = T_x (X)$, а $T_x (X)$ --- касательное подпространство


к $X$ в точке $x$, рассматриваемое как $n$-мерное


проективное пространство $\Bbb{P}^n$.


Число $r$ называют также


{\em рангом} многообразия $X$, $r = \rank X$.


Случай $r= 0 $ является тривиальным: многообразие $X$ оказывается


$n$-плоскостью.



Пусть $X \subset \Bbb{P}^N$ --- почти всюду гладкое


$n$-мерное многообразие с вырожденным гауссовым отображением.


Предположим, что $0 < \rank \gamma = r < n$.


Обозначим через $L$ слой отображения


$$


\gamma:L = \gamma^{-1} (T_x) \subset X;\quad \dim L = n - r = l.


$$



Как доказано в (\cite{AG04}, теоремa 3.1, с.\,95),


многообразие с вырожденным гауссовым отображением ранга


$r$ расслаивается на плоские слои


$L$ размерности $l$, вдоль которых касательное подпространство


$T_x (X)$ постоянно.


 Касательное подпространство $T_x (X)$ является постоянным, когда


точка $x$ пробегает множество регулярных точек слоя $L$.


По этой причине мы обозначаем его $T_L$, $L \subset T_L$.


Пара $(L, T_L)$ на многообразии $X$ зависит от $r$ параметров.



Вышеописанное слоение на многообразии $X$


называется {\em слоением Монжа--Ампера}.





Многообразия ранга $r < n$ являются многомерными аналогами


развертывающихся поверхностей трехмерного евклидова пространства.





Изложение основных результатов геометрии многообразий с


вырожденным гауссовым отображением и дальнейшую литературу по этому


вопросу можно найти в (\cite{AG93}, гл.\,4), а также в~\cite{AG04}.





Данный параграф посвящен отысканию основных уравнений


многообразия $X$ с вырожденным гауссовым отображением


размерности $n$ и ранга $r$


в проективном пространстве $\Bbb{P}^N$.





В дальнейшем используемые индексы будут пробегать следующие области


значений:


$$
a, b, c = 1, \dotsc, l; \ \ p, q = l + 1, \dotsc , n; \ \ \alpha,
\beta = n + 1, \dotsc , N.
$$





Точка $x \in X$ называется {\it регулярной точкой} отображения


$\gamma$ (или многообразия $X$), если $\dim T_x (X) = \dim X = n$.





Свяжем с многообразием $X$ семейство подвижных реперов


$\{A_u\}$, $u = 0, 1,\dotsc , N$, таким образом, что точка $A_0 = x$


является регулярной точкой слоя $L$ многообразия $X$, точки $A_a$


принадлежат слою $L$ слоения Монжа--Ампера, проходящему через точку


$A_0$, точки $A_p$ вместе с точками $A_0, A_a$ определяют касательное


подпространство $T_L (X)$ к многообразию $X$, а точки $A_\alpha$


расположены вне подпространства $T_L (X)$.





Уравнения инфинитезимального перемещения подвижного репера


$\{A_u\}$ имеют вид


\begin{equation*}


%\label{eq:1}


 dA_u = \omega_u^v A_v, \quad u, v = 0, 1, \dotsc , N,
\end{equation*}
где


$\omega_u^v$ --- 1-формы, удовлетворяющие структурным уравнениям


проективного пространства $\Bbb{P}^N$:


\begin{equation*}


%\label{eq:2}


d \omega_u^v = \omega_u^w \wedge


\omega_w^v, \quad u, v, w = 0, 1, \dotsc , N.


\end{equation*}



В результате вышеуказанной специализации подвижного репера получим


следующие уравнения многообразия $X$ (\cite{AG04}, секция 3.1):


\begin{gather}


% \label{eq:3}


\omega_0^\alpha = 0,\quad 
% \begin{equation}\label{eq:4}
\omega_a^\alpha = 0,\notag \\
%
\label{eq:5}
\omega_p^\alpha = b_{pq}^\alpha \omega^q, \quad b_{pq}^\alpha =
b_{qp}^\alpha, \\
%
\label{eq:6}
\omega^p_a = c^p_{aq} \omega^q.
\end{gather}



В этих уравнениях 1-формы $\omega^q := \omega_0^q$


являются базисными формами гауссова образа $\gamma (X)$ многообразия $X$,


а величины $b_{pq}^\alpha$ образуют второй фундаментальный тензор


многообразия $X$ в точке $x = A_0$.


Величины $b_{pq}^\alpha$ и $c^p_{aq}$ связаны следующими соотношениями:


\begin{equation}


\label{eq:7}


b_{sq}^\alpha c^s_{ap} = b_{sp}^\alpha c^s_{aq}.
\end{equation}



Уравнения (1) и (2) называются {\em основными уравнениями}


многообразия $X$ с вырожденным гауссовым отображением


(\cite{AG04}, cекция 3.1).





Отметим, что при преобразованиях точек $A_p$ величины


$c^p_{aq}$ преобразуются как тензоры по отношению к индексам


$p$ и $q$. По отношению к индексу $a$, величины


$c^p_{aq}$ тензора не образуют.


Тем не менее, при преобразованиях точек $A_0$ и $A_a$ величины


$c^p_{aq}$ вместе с единичным тензором $\delta_q^p$


преобразуются как тензоры.


По этой причине система величин $c^p_{aq}$ называется


{\em квазитензором}.





Обозначим через $B^\alpha$ и $C_a$ следующие $r \times r$-матрицы,


состоящие из коэффициентов уравнений (1) и (2):


$$


B^\alpha = (b^\alpha_{pq}),\quad C_a = ( c^p_{aq}).


$$
В некоторых случаях мы будем использовать единичную матрицу


$C_0 = (\delta^p_q)$ и индекс


$i = 0, 1, \dotsc, l$, т.\,е. $\{i\} = \{0, a\}$.


Соотношения (1) и (3) можно объединить


и записать в матричном виде


$$


(B^\alpha C_i)^T = (B^\alpha C_i),


$$
что означает симметричность матриц


$$


H^\alpha_i = B^\alpha C_i = (b^\alpha_{qs} c^s_{ip}).


$$





Квадратичные формы


$$


\Phi^\alpha = b_{pq}^\alpha \omega^p \omega^q


$$
являются {\it вторыми фундаментальными формами} многообразия


$X$ в точке $x = A_0$, а формы


$$
\Phi^\alpha = b^\alpha_{ps} (\delta^s_q x^0 + c^s_{aq} x^a) 
\omega^p \, \omega^q
$$
являются {\it вторыми фундаментальными формами} многообразия $X$


в точке $x = x^0 A_0 + x^a A_a \in L$.



Пусть $\{\alpha^u\}$ --- корепер (или тангенциальный репер)


в пространстве $(\Bbb{P}^{N})^*$,


дуальный к реперу $\{A_u\}$.


Тогда гиперплоскости $\alpha^u$ репера $\{\alpha^u\}$ связаны с


вершинами репера $\{ A_u\}$ условиями


\begin{equation}


\label{eq:9}


(\alpha^u, A_v) = \delta^u_v.
\end{equation}
Условия (4) означают, что гиперплоскость $\alpha^u$ содержит все


точки $A_v$, $v \neq u$, а также что выполнено следующее условие


нормировки: $(\alpha^u, A_u) = 1$.



Уравнение


$$\xi_\beta \alpha^\beta = 0


$$
определяет систему касательных гиперплоскостей, проходящих через


касательное подпространство $T_L (X)$.


Эта система касательных гиперплоскостей определяет систему вторых


фундаментальных форм


$$


\mathrm{II} = \xi_\beta b_{pq}^\beta \omega^p \omega^q


$$
и систему вторых фундаментальных тензоров


$\xi_\beta b_{pq}^\beta$
многообразия $X$ в точке $x = A_0$.





Многообразие $X$ с вырожденным гауссовым отображением является


{\it дуально невырожденным}, если размерность его дуального


многообразия $X^*$ равна $N - l -1$.


 Согласно обобщенной теореме Гриффитса--Харриса (теорема 3.2 и


следствие 3.3 в \cite{AG04}, с.\,97--99) {\it многообразие


$X$ с вырожденным гауссовым отображением является дуально


невырожденным тогда и только тогда, когда


в любой гладкой точке $x \in X$ найдется по крайней мере одна


невырожденная вторая фундаментальная форма в системе


вторых фундаментальных форм $\xi_\alpha b_{pq}^\alpha \omega^p


\omega^q$ многообразия~$X$.}





В дальнейшем будем рассматривать только дуально невырожденные


многообразия $X$ с вырожденным гауссовым отображением.








\section{Фокальные образы многообразия\protect\\ с вырожденным гауссовым


отображением}



Предположим, что $X$ --- многообразие с вырожденным гауссовым


отображением размерности $n$ и ранга $r$ в пространстве


$\Bbb{C} \Bbb{P}^N$.


 Как было ранее замечено, такое многообразие несет на себе


$r$-параметрическое семейство $l$-мерных плоскостных образующих $L$


размерности $l = n - r$.


Пусть $x = x^0 A_0 + x^a A_a$ --- произвольная точка образующей $L$.


Для такой точки имеем


$$
d x = (dx^0 + x^0 \omega_0^0 + x^a \omega_a^0) A_0 + (dx^a + x^0
\omega^a + x^b \omega_b^a) A_a + (x^0 \omega^p + x^a \omega_a^p)A_p.
$$
В соответствии с (2) отсюда следует


\begin{equation}


\label{eq:11}


d x \equiv (x^0 \delta_q^p + x^a c^p_{aq}) A_p \omega^q \pmod{L}.
\end{equation}
Матрица $(J^p_q) = (x^0 \delta_q^p + x^a c^p_{aq})$ является


{\em матрицей Якоби} отображения


$\gamma: X \rightarrow \Bbb{G} (n, N)$, а опре\-делитель


$$


J (x) = \det (J^p_q )


$$
этой матрицы является {\em якобианом} отображения $\gamma$.


Легко видеть, что $J (x) \neq 0$ в регулярных точках и $J(x) = 0$


в особых точках.



В соответствии с (2) и (5) множество особых точек образующей $L$


многообразия $X$ определяется уравнением


\begin{equation}


\label{eq:12}


\det (\delta_q^p x^0 + c^p_{aq} x^a) = 0.
\end{equation}
 Таким образом, это множество является


{\em алгебраической гиперповерхностью размерности $l- 1$


и степени $r$ в образующей $L$}.


Эта гиперповерхность (в $L$) называется


{\em фокусной гиперповерхностью} и обозначается $F_L$.





Поскольку при $x^a = 0$ левая часть уравнения (6) принимает вид


$$


\det (x^0 \delta^p_q) = (x^0)^r,


$$
то отсюда следует, что точка $A_0$ является регулярной точкой образующей $L$.






Будем называть касательную гиперплоскость


$\xi = (\xi_\alpha)$ {\em особой} (или {\em фокусной


гиперплоскостью}), если


\begin{equation}


\label{eq:13}


\det (\xi_\alpha b^\alpha_{pq}) = 0,
\end{equation}
т.\,е. если матрица $(\xi_\alpha b^\alpha_{pq})$ имеет пониженный ранг.


Условие (7) представляет собой уравнение степени $r$ по отношению к


тангенциальным координатам $\xi_\alpha$ гиперплоскости


$\xi$, содержащей касательное подпространство $T_L (X)$.



Поскольку в данной работе рассматриваются только дуально


невырожденные многообразия с вырожденным гауссовым отображением, то


{\it найдется по крайней мере одна невырожденная форма в системе


вторых фундаментальных форм многообразия $X$} (см. конец \S\,1).


Отсюда следует, что уравнение (7) определяет в пространстве


$\Bbb{P}^N$ алгебраический гиперконус степени $r$,


вершиной которого является касательное подпространство $T_L (X)$.


Этот гиперконус называется {\em фокусным гиперконусом} и


обозначается $\Phi_L$ (\cite{AG93}, с.\,119).





Заметим, что в случае, когда многообразие $X$ является дуально


вырожденным, уравнения (7) на этом многообразии удовлетворяются


тождественно и многообразие $X$ не имеет фокусных гиперконусов.



Фокусная гиперповерхность $F_L \subset L$ и фокусный гиперконус


$\Phi_L$ с вершиной $T_L$ называются {\em фокальными образами}


многообразия $X$ с вырожденным гауссовым отображением.






\section{Конгруэнции и псевдоконгруэнции в проективном пространстве}



Рассмотрим семейство $Y$ \ $l$-мерных подпространств


$L$, $\dim L = l$, в проективном пространстве $\Bbb{P}^n$,


зависящее от $r = n - l$ параметров.


 Будем предполагать, что через всякую точку $x \in \Bbb{P}^n$


проходит не более конечного числа подпространств $L$


этого семейства.


 Если ограничиться рассмотрением малой окрестности подпространства


$L$, то можно считать, что только одно подпространство


$L \in Y$ проходит через произвольную точку $x \in L$.


 Семейства такого вида в пространстве $\Bbb{P}^n$ называются


{\it конгруэнциями}.





Дуальным образом для конгруэнции $Y$ $l$-мерных подпространств


в $\Bbb{P}^n$ является {\it псевдоконгруэнция} $Y^*$,


представляющая собой $r$-параметрическое семейство подпространств


размерности $r - 1$.


 Всякая гиперплоскость $\xi \subset \Bbb{P}^n$ содержит не более


чем конечное число подпространств $L^* \in Y^*$.


 Однако, если рассматривать достаточно малую окрестность


подпространства $L^*$ псевдоконгруэнции $Y^*$,


то гиперплоскость $\xi$ будет содержать только единственное


подпространство~$L^*$.





 Установим связь теории многообразий с


вырожденным гауссовым отображением в проективных пространствах с


теорией конгруэнций и псевдоконгруэнций подпространств и покажем,


как эти две теории могут быть применены к построению


индуцированных связностей на подмногообразиях проективных пространств,


а также других пространств, снабженных проективной структурой.



Итак, рассмотрим в проективном пространстве $\Bbb{P}^n$


конгруэнцию $Y$ \ $l$-мерных подпространств $L$.


Свяжем с любым элементом $L$ конгруэнции $Y$ семейство


проективных реперов


$\{A_0, A_1, \dotsc , A_n\}$, выбранных таким образом, что точки


$A_0, A_1, \dotsc , A_l$ принадлежат подпространству $L$,


а точки $A_{l+1}, \dotsc , A_n$ расположены вне этого подпространства.


Уравнения инфинитезимального перемещения для таких реперов имеют вид


\begin{equation}


\label{14}


\begin{cases}


dA_i = \omega_i^j A_j + \omega_i^p A_p, \\
dA_p = \omega_p^i A_i + \omega_p^q A_q, 
\end{cases}


\end{equation}
где $i, j = 0, 1, \dotsc, l$; $p, q = l + 1, \dotsc , n$, а


$L = A_0 \wedge A_1\wedge \dotsb \wedge A_l$ --- образующая рассматриваемой


конгруэнции $Y$.


 Поскольку эта образующая зависит от $r$ параметров и является


фиксированной при $\omega_i^p = 0$, то формы


$\omega_i^p$ линейно выражаются через дифференциалы этих


$r$ параметров или через линейно независимые 1-формы


$\theta^p$, являющиеся линейными комбинациями этих дифференциалов:


\begin{equation}


\label{15}


 \omega_i^p = c_{iq}^p \theta^q.
\end{equation}
 При допустимых линейных преобразованиях базисных форм


$\theta^p$ матрицы $C_i = (c_{iq}^p)$ преобразуются по тензорному


закону по отношению к индексам $p$ и $q$.





Точка $F \in L \subset Y$ называется {\it фокусом} образующей


$L$, если $dF \in L$ при некотором условии на базисные формы $\theta^p$.


Для отыскания фокусов представим их в виде $F = x^i A_i$.


Тогда


$$
dF \equiv x^i \omega_i^p A_p \pmod{L};


$$
поэтому фокусы определяются системой уравнений


$$


x^i \omega_i^p = 0.


$$
В соответствии с (9) эта система принимает вид


\begin{equation}\label{16}
x^i c_{iq}^p \theta^q = 0.
\end{equation}


Система (10) имеет нетривиальное решение по отношению к


формам $\theta^q$ тогда и только тогда, когда


\begin{equation}\label{17}
 \det (x^i c_{iq}^p) = 0.
\end{equation}
Уравнение (11) определяет на $L$ фокусную гиперповерхность $F_L$,


представляющую собой алгебраическую гиперповерхность степени $r$.





Предположим, что точка $A_0$ рассматриваемого подвижного репера не


принадлежит гиперповерхности $F_L$.


Тогда 1-формы $\omega_0^p$ являются линейно независимыми, и мы можем


взять эти формы в качестве базисных форм конгруэнции $Y$.


 В результате уравнения (9) приобретают вид


\begin{equation}\label{18}
 \omega_a^p = c_{aq}^p \omega_0^q,
\end{equation}
где $a = 1, \dotsc , l$, а $c_{0q}^p = \delta_q^p$.


Теперь уравнения (12) совпадают с уравнениями (2).


Поэтому уравнение (11) фокусной гиперповерхности $F_L$


запишется в форме


\begin{equation}\label{19}
 \det (x^0 \delta_q^p + x^a c_{aq}^p) = 0.
\end{equation}
Уравнение (13) совпадает с уравнением (6), определяющим фокусы на


плоской образующей $L$ многообразия $X$ с вырожденным гауссовым


отображением ранга $r$.


Однако в отличие от \S\,1 эти величины


$c_{aq}^p$ не связаны никакими соотношениями типа (3), поскольку в


данном случае нет матриц $B^\alpha = (b_{pq}^\alpha)$.


Таким образом, фокусные гиперповерхности $F_L$, определяемые


уравнением (13), являются произвольными детерминантными многообразиями


на образующих $L$ рассматриваемой конгруэнции $Y$.





В частности, если $l = 1$ и $n = r + 1$, то $Y$ становится


прямолинейной конгруэнцией.


Уравнение (13), определяющее фокусные гиперповерхности


$F_L$ такой конгруэнции, принимает вид


\begin{equation*}


%\label{20}
 \det (x^0 \delta_q^p + x^1 c_{1q}^p) = 0.
\end{equation*}


 Следовательно, каждая из фокусных гиперповерхностей $F_L$ конгруэнции


$Y$ распадается на $r$ вещественных или комплексных точек при


условии, что каждая из них считается столько раз, какова ее кратность.


 Каждая из этих точек описывает


{\it фокальное многообразие} в пространстве $\Bbb{P}^n$,


касательное к образующим $L$ конгруэнции $Y$.






Рассмотрим далее псевдоконгруэнцию $Y^*$ в пространстве $\Bbb{P}^n$.


Ее образующая $L^*$ имеет размерность $r - 1$ и зависит от


$r$ параметров.


 Поместим точки $A_p$, $p= l + 1, \dotsc , n$, $l = n - r$,


рассматриваемого подвижного репера


на образующую $L^* \in Y^*$, а точки $A_i$, $ i = 0, 1, \dotsc , l$,


вне образующей $L^*$.


Уравнения инфинитезимального перемещения таких реперов снова имеют


вид (8), однако в данном случае 1-формы $\omega_p^i$


являются линейными комбинациями базисных форм


$\theta^p$, определяющих смещение образующей $L^* = A_{l+1}


\wedge \dotsb \wedge A_n$.


Таким образом, в этом случае имеем


$$


\omega_p^i = b_{pq}^i \theta^q


$$
и


\begin{equation}\label{22}
dA_p = \omega_p^q A_q + b_{pq}^i \theta^q A_i.
\end{equation}

Рассмотрим гиперплоскость $\xi$, проходящую через образующую


$L^* \in Y^*$.


 По отношению к рассматриваемому подвижному реперу уравнение


гиперплоскости $\xi$ имеет вид $\xi_i x^i = 0$, где $\xi_i$~---


тангенциальные координаты этой гиперплоскости.


 Гиперплоскость $\xi$, которая кроме образующей


$L^*$ содержит также бесконечно близкую образующую ${}'L^*$,


определяемую точками $A_p$ и $dA_p$, называется


{\it фокусной гиперплоскостью}.


В соответствии с (14) уравнения, определяющие фокусные


гиперплоскости, имеют вид


\begin{equation}\label{23}


\xi_i b_{pq}^i \theta^q = 0.
\end{equation}
Система уравнений (14) определяет смещение образующей $L^*$


тогда и только тогда, когда система (15) имеет нетривиальное решение по


отношению к формам $\theta^q$.


Необходимым и достаточным условием существования такого


нетривиального решения является обращение в нуль определителя системы (15):


\begin{equation}\label{24}
\det (\xi_i b_{pq}^i) = 0.
\end{equation}



Уравнение (16) определяет семейство фокусных гиперплоскостей,


проходящих через образующую $L^* \in Y^*$.


 Это семейство представляет собой алгебраический гиперконус степени


$r$, вершиной которого является образующая $L^*$.


 Заметим, что уравнение (16) аналогично уравнению (7) фокусного


гиперконуса $\Phi_L$ многообразия с вырожденным гауссовым


отображением.






\section{Нормализованные многообразия в многомерном проективном


пространстве}





{\bf 1.} Рассмотрим гладкое $r$-мерное многообразие $X$


в проективном пространстве $\Bbb{P}^n$, $r < n$.


Дифференциальная геометрия такого многообразия


беднее, чем дифференциальная геометрия многообразия в евклидовом


пространстве или пространствах постоянной кривизны


$\Bbb{S}^n$ и $\Bbb{H}^n$, где $\Bbb{S}^n$ и


$\Bbb{H}^n$ означают соответственно $n$-мерные эллиптическое и


гиперболическое пространства.


 С окрестностью первого порядка точки


$x \in X \subset \Bbb{P}^n$ ассоциируется только касательное


подпространство $T_x (X)$.


Например, в (\cite{AG04}, секция 1.4)  показано, что для того чтобы


обогатить дифференциальную геометрию кривой на проективной плоскости


$\Bbb{P}^2$, необходимо использовать дифференциальные продолжения


уравнения кривой достаточно высоких порядков.





Однако мы можем обогатить дифференциальную геометрию многообразия


$X\subset\Bbb{P}^n$, снабжая $X$ дополнительной конструкцией,


состоящей из подпространства $N_x (X)$ размерности $n - r$ такого, что


$T_x (X) \cap N_x (X) = x$, и $(r - 1)$-мерного подпространства


$K_x (X)$, $K_x (X) \subset T_x (X)$, $x \notin K_x (X)$.


 Будем обозначать эти подпространства символами $N_x$ и $K_x$ и


называть {\it нормалями первого и второго рода} (или просто


{\it первой} и {\it второй нормалями}) многообразия $X$


соответственно (\cite{N76}, с.\,198).





Семейство первых нормалей образует


{\it конгруэнцию} $N$, а семейство вторых нормалей образует


{\it псевдоконгруэнцию} $K$ в пространстве $\Bbb{P}^n$.


Если с каждой точкой $x \in X$ ассоциированы единственная первая


нормаль $N_x$ и единственная вторая нормаль $K_x$, то многообразие


$X$ называется {\it нормализованным} (\cite{N76}, с.\,198;


\cite{AG93}, гл.\,6).



Как будет видно из дальнейшего, в случае многообразий в евклидовом


пространстве $\Bbb{E}^n$, а также неевклидовых пространствах


$\Bbb{S}^n$ и


$\Bbb{H}^n$, первые и вторые нормали определяются геометрией этих


пространств, в то время как в случае многообразий в аффинном


пространстве $\Bbb{A}^n$ и в проективном пространстве


$\Bbb{P}^n$ эти нормали должны быть выбраны искусственно


или же для нахождения их следует использовать окрестности точки $x \in X$


высоких порядков.


 В этой статье мы будем применять первый метод.


 Заметим, что второй метод сопряжен со значительными вычислительными


трудностями.


 Подробное изложение этого метода и ссылки на соответствующую литературу


можно найти в (\cite{AG93}, гл.\,6, 7; \cite{N76}, гл.\,5).





 Итак, рассмотрим нормализованное многообразие


$X$ размерности $r$ в проективном пространстве $\Bbb{P}^n$.


 Свяжем с $X$ семейство проективных реперов $\{A_0, A_1, \dotsc ,


A_n\}$ таким образом, что $A_0 = x$, $A_a \in N_x$, $a = 1, \dotsc,


l$, где $l = n - r$, и $A_p \in K_x$, $ p = l + 1, \dotsc, n$.


Уравнения инфинитезимального перемещения таких реперов имеют вид


\begin{equation}\label{eq:25}
\begin{cases}
dA_0 = \omega_0^0 A_0 \phantom{+ \omega_a^b A_b.}


\ + \omega^p A_p , \\
dA_a = \omega_a^0 A_0 + \omega_a^b A_b + \omega^p_a A_p, \\
dA_p = \omega_p^0 A_0 + \omega^a_p A_a + \omega_p^q A_q.


\end{cases}
\end{equation}



Уравнения (17) показывают, что рассматриваемое семейство подвижных


реперов удовлетворяет системе дифференциальных уравнений


\begin{equation}\label{eq:26}
\omega^a = 0,


\end{equation}
а 1-формы $\omega^p$ являются базисными формами, поскольку они


определяют смещение точки $A_0 = x$ вдоль многообразия $X$.


 Внешнее дифференцирование уравнения (18) и использование леммы


Картана приводит к уравнениям


\begin{equation}\label{eq:27}
\omega_p^a = b_{pq}^a \omega^q, \;\; b_{pq}^a = b_{qp}^a.


\end{equation}


 Величины $b^a_{pq}$ образуют тензор и являются коэффициентами вторых


фундаментальных форм многообразия


$X$ в точке $x$ (\cite{AG04}, секция 2.1):


$$


\Phi^a = b_{pq}^a \omega^p \omega^q.


$$

{\bf 2.}


Точки $A_p$ принадлежат касательному подпространству $T_x(X)$.


Будем предполагать, что эти точки принадлежат второй нормали


$K_x\subset T_x (X)$, $K_x = A_{l+1} \wedge \dotsb \wedge A_n$.


Тогда при $\omega^p = 0$ 1-формы $\omega_p^0$ также должны


обращаться в нуль, в результате чего имеем


\begin{equation}\label{eq:29}
\omega_p^0 = l_{pq} \omega^q.
\end{equation}



Далее, поместим точки $A_a$ рассматриваемого подвижного репера


на первую нормаль $N_x$ многообразия $X$,


$N_x = A_0 \wedge A_{1} \wedge \dotsb \wedge A_l$.


Тогда при $\omega^p = 0$ получаем $\omega_a^p = 0$, откуда


\begin{equation}\label{eq:30}
\omega_a^p = c_{a q}^p \omega^q.
\end{equation}



 Рассмотрим точку $y \in N_x$ на первой нормали.


 Имеем $y = y^0 A_0 + y^a A_a$.


 Дифференцируя это соотношение и используя (17), получаем


\begin{equation}\label{eq:31}


 dy = (dy^0 + y^0 \omega_0^0 + y^a
\omega_a^0) A_0 + (y^0 \omega^p + y^a \omega_a^p) A_p + (dy^a +
y^b \omega_b^a) A_a.
\end{equation}

Точка $y$ является {\it фокусом} первой нормали $N_x$, если


$dy \in N_x$.


 В соответствии с (22) из этого условия следует


$$
y^0 \omega^p + y^a \omega_a^p = 0.
$$


Применяя соотношения (21), находим


$$


(y^0 \delta_q^p + y^a c_{a q}^p) \omega^q = 0.


$$
Эта система имеет нетривиальное решение по отношению к формам


$\omega^q$ тогда и только тогда, когда


\begin{equation}\label{eq:32}
\det (y^0 \delta_q^p + y^a c_{a q}^p) = 0.
\end{equation}
 Уравнение (23) отличается от уравнения (13) только обозначениями 


и определяет фокусную гиперповерхность $F_x$ в образующей $N_x$


конгруэнции первых нормалей, ассоциированной с многообразием $X$.


 Из уравнения (23) следует, что точка $x \in X$ с координатами


$x^0 = 1$, $x^a = 0$ не принадлежит фокусной гиперповерхности $F_x$.



Найдем фокусные гиперконусы $\Phi_x$ псевдоконгруэнции


$K$ вторых нормалей многообразия $X$.


 Гиперконус $\Phi_x$ образован гиперплоскостями


$\xi$ пространства $\Bbb{P}^n$, содержащими


вторую нормаль $K_x = A_{l+1} \wedge \dotsb \wedge A_n \subset


T_x (X)$ и бесконечно близкую нормаль $K_x + d K_x$, содержащую не


только точки $A_p$, но также и точки


$$


d A_p \equiv \omega_p^0 A_0 + \omega_p^a A_a \pmod{N_x}.


$$


Таким образом, тангенциальные координаты $\xi_0$ и $\xi_a$


такой гиперплоскости удовлетворяют уравнениям


$$


\xi_0 \omega_p^0 + \xi_a\omega_p^a = 0.


$$
В соответствии с (20) и (19) из этого уравнения следует


$$


(\xi_0 l_{pq} + \xi_a b^a_{pq})\, \omega^q = 0.


$$
Эта система имеет нетривиальное решение по отношению к формам


$\omega^q$ тогда и только тогда, когда обращается в нуль ее


определитель:


\begin{equation}\label{eq:33}
\det (\xi_0 l_{pq} + \xi_a b^a_{pq}) = 0.
\end{equation}
Уравнение (24) определяет алгебраический гиперконус порядка $r$, вершиной


которого является образующая $K_x$ псевдоконгруэнции $K$ вторых нормалей.


 Этот гиперконус называется {\it фокусным гиперконусом}


псевдоконгруэнции $K$.


\medskip





{\bf 3.}


Рассмотрим теперь касательное и нормальное расслоения,


ассоциированные с нормализованным многообразием $X$.


Базой каждого из этих расслоений является само многообразие $X$,


слоями касательного расслоения являются касательные подпространства


$T_x$, а слоями нормального расслоения являются первые нормали $N_x$.





Пусть ${}'x = x + x^p A_p$ --- произвольная точка касательного


подпространства $T_x$, а $\boldsymbol{x} = {}'x - x =


 x^p A_p$ --- вектор в касательном расслоении $T(X)$.


 Дифференциал этого вектора имеет вид


\begin{equation}\label{eq:34}


 d\boldsymbol{x} =


 (dx^p + x^q \omega_q^p) A_p + x^p (l_{pq} A_0 + b^a_{pq} A_a) \omega^q.


\end{equation}


Первое слагаемое в правой части соотношения (25) принадлежит


касательному подпространству $T_x$, а второе слагаемое принадлежит


нормали $N_x$. \ 1-форма


$$


D x^p = dx^p + x^q \omega_q^p


$$
называется {\it ковариантным дифференциалом} вектора


$\boldsymbol{x} = (x^p)$ по отношению к аффинной связности $\gamma_t$,


индуцированной на многообразии $X$ нормализацией $(N, K)$.


 1-формы $\omega_q^p$ суть компоненты {\it формы связности}


$\omega =\{\omega_q^p\}$ связности $\gamma_t$.





Векторное поле $\boldsymbol{x}$ называется {\it параллельным}


в связности $\gamma_t$, если его ковариантный дифференциал


$Dx^p$ обращается в нуль, т.\,е. если


$$


Dx^p = dx^p + x^q \omega_q^p = 0.


$$

Найдем внешние дифференциалы компонент $\omega_q^p$


формы связности $\omega$.


Из формул (19), (20) и (21) следует, что эти внешние


дифференциалы имеют вид


\begin{equation}\label{eq:36}


d \omega_q^p = \omega_q^s \wedge \omega_s^p + (l_{qs} \delta_t^p +
b^a_{qs} c_{at}^p) \omega^s \wedge \omega^t.
\end{equation}
2-форма


$$


\Omega_q^p = d \omega_q^p - \omega_q^s \wedge \omega_s^p


$$
называется {\it формой кривизны} аффинной связности


$\gamma_t$, индуцированной на многообразии $X$.


Из уравнения (26) следует, что


$$


\Omega_q^p = \tfrac{1}{2} R^p_{qst} \omega^s \wedge \omega^t,


$$
где


\begin{equation}\label{eq:38}
R^p_{qst} = l_{qs} \delta_t^p + b^a_{qs} c_{at}^p - l_{qt}
\delta_s^p - b^a_{qt} c_{as}^p
\end{equation}
--- {\it тензор кривизны} аффинной связности $\gamma_t$ на $X$.


Уравнения (27) позволяют вычислять тензоры кривизны для различных


нормализаций многообразия $X$.



Если $R^p_{qst} = 0$ на многообразии $X$, то аффинная связность


$\gamma_t$ на $X$ является {\it плоской},


и параллельное перенесение вектора $\boldsymbol{x}$


не зависит от пути (см., напр., \cite{N76}, с.\,118; \cite{KN}, с.\,70).


\medskip





{\bf 4.}


Рассмотрим далее векторное поле $\boldsymbol{y}$


в нормальном расслоении $N (X)$.


Вектор $\boldsymbol{y}$ определяется точкой $x$


и точкой $y = y^0 A_0 + y^a A_a$ слоя $N_x \subset N (X)$.


Дифференциал точки $y$ определяется уравнением (22).





1-форма


$$


D y^a = dy^a + y^b \omega_b^a


$$
называется {\it ковариантным дифференциалом}


векторного поля $\boldsymbol{y}$ в нормальном расслоении $N (X)$,


а формы $\omega_a^b$ являются компонентами


{\it формы связности нормальной связности}


$\gamma_n$ на нормализованном многообразии $X$


(см., напр., \cite{Ca01}, с.\,242; дополнительную информацию


о нормальной связности можно найти в \cite{AG95} и


\cite{AG93}, cекция 6.3). \ 2-форма


$$


\Omega_b^a = d \omega_b^a - \omega_b^c \wedge \omega_c^a


$$
называется {\it формой кривизны} нормальной связности $\gamma_n$.


 Отметим, что Картан в \cite{Ca01} называл эту форму


{\it гауссовым кручением} вложенного многообразия $X$.



Дифференцируя формы $\omega_b^a$ и применяя формулы


(19) и (20), для формы кривизны  находим выражение


$$


\Omega_b^a = \tfrac{1}{2}R^a_{bst} \omega^s \wedge \omega^t,


$$


где


\begin{equation}\label{eq:41}
R^a_{bst} = c_{bs}^p b^a_{pt}- c_{bt}^p b^a_{ps}.
\end{equation}
 Тензор $R^a_{bst}$ называется {\it тензором нормальной кривизны}


многообразия $X$.



Вторые нормали $K_x$, ассоциированные с многообразием $X$, позволяют


получить распределение $\Delta_y$ $r$-мерных подпространств,


ассоциированное с $X$.


 Элементы этого распределения $\Delta_y$


являются линейными оболочками точек $y \in N_x$ и вторых нормалей


$K_x$, $\Delta_y = y \wedge K_x$.


Из (22) следует, что распределение


$\Delta_y$ определяется системой уравнений


\begin{equation}\label{eq:42}
dy^a + y^b \omega^a_b = 0.
\end{equation}
В общем случае система уравнений (29) не является вполне


интегрируемой, и когда точка $x$ движется вдоль замкнутого контура


$l \subset X$, соответствующая точка


$y$ не описывает замкнутого контура.



Однако точка $y$ будет описывать замкнутый контур $l'$, если система (29)


является вполне интегрируемой.


 Условием полной интегрируемости системы (29) является обращение в


нуль тензора кривизны (28) нормальной связности многообразия $X$.


 В этом случае распределение $\Delta_y$, определенное системой


(29), является вполне интегрируемым, и замкнутые контуры


$l'$ лежат на интегральных многообразиях этого распределения.


Эти интегральные многообразия составляют


$(n - r)$-параметрическое семейство $r$-мерных подмногообразий


$X (y)$, которые ``параллельны'' многообразию


$X$ в том смысле, что подпространства $T_x (X)$ и $T_x (X (y))$


проходят через одну и ту же вторую нормаль $K_x$.


\medskip



{\bf 5.}


Нормализация многообразия $X \subset \Bbb{P}^n$


называется {\it центральной}, если все первые нормали


$N_x$ этой нормализации образуют связку с $(l-1)$-мерной вершиной $S$.




\begin{thm}
Нормализация нормализованного многообразия


$X \subset \Bbb{P}^n$ является центральной тогда и только тогда,


когда величины $l_{pq}$ и $c^p_{aq}$ в уравнениях $(20)$ и $(21)$


обращаются в нуль\/${:}$


\begin{equation}
l_{pq} = 0, \quad c^p_{aq} = 0.
\end{equation}
\end{thm}

\begin{pf}


{\bf Необходимость.}


 Предположим, что нормализация многообразия $X \subset


\Bbb{P}^n$ является центральной нормализацией


с $(l-1)$-мерной вершиной $S$.


 Поместим точки $A_a$ в эту вершину $S$.


Тогда


$$


d A_a = \omega_a^b A_b.


$$
При этом из уравнения (17) следует


\begin{equation}
\omega_a^0 = 0, \;\; \omega_a^p = 0.
\end{equation}
В соответствии с соотношениями (20) и (21) это означает (30).


\smallskip



{\bf Достаточность.}


Из (30) следует (31), откуда $d A_a = \omega_a^b A_b$.


Таким образом, подпространство $S = A_{1} \wedge \dotsb \wedge A_l$


является постоянным.


 Все $l$-мерные первые нормали $N_x$ проходят


через $S$, и нормализация многообразия $X$ является центральной с


$(l-1)$-мерной вершиной $S$.


\end{pf}



\begin{cornonum}
Индуцированная аффинная связность $\gamma_t$ и нормальная связность


$\gamma_n$ центрально нормализованного многообразия $X \subset


\Bbb{P}^n$ являются плоскими.


\end{cornonum}

\begin{pf} Поскольку для центрально нормализованного многообразия $X$


выполняются соотношения (30),


а тензор кривизны индуцированной аффинной связности


$\gamma_t$ имеет вид (27), то


$$


R^p_{qst} = 0,


$$
т.\,е. тензор кривизны индуцированной аффинной связности $\gamma_t$ центрально


нормализованного многообразия обращается в нуль.





Аналогичным образом из формулы (28) следует, что тензор кривизны


нормальной связности $\gamma_n$ центрально нормализованного


многообразия также обращается в нуль.





 Отметим, что оба этих факта вытекают также из того, что центрально


нормализованное многообразие $X \subset \Bbb{P}^n$ может быть


биективно спроектировано на


$r$-мерное подпространство $T$, которое является дополнительным к


вершине $S$ расслоения первых нормалей


$N_x$, и геометрия многообразия $X$, индуцированная этой центральной


нормализацией, эквивалентна плоской геометрии пространства $T$.


\end{pf}

В \cite{A52} были найдены необходимые и достаточные


условия, при которых нормализация многообразия


$X$ в аффинном пространстве $\Bbb{A}^{N}$ является центральной или


тривиальной.


Тривиальной нормализацией многообразия $X$ в аффинном пространстве


$\Bbb{A}^N$ называется нормализация, для которой все первые нормали


$N_x$ параллельны некоторой постоянной $l$-плоскости


(т.е. образуют пучок параллельных $l$-плоскостей).





В обозначениях, принятых в данной работе,


условия центральности нормализации,


полученные в \cite{A52}, имеют вид


$$


c_{aq}^p = \delta_q^p c_a,


$$
где $\delta_q^p$ --- символы Кронекера, $c_a$ --- $(1, 0)$-тензоры,


а условия, при которых нормализация является тривиальной, имеют вид


$$


c_{aq}^p = 0.


$$
Однако в случае аффинного пространства (и, в частности, в случае


евклидова пространства) всегда выполняется $l_{pq} = 0$.


 Кроме того, в проективном случае (так же, как и в аффинном случае),


тривиальная нормализация является центральной нормализацией, вершина


$S$ которой принадлежит бесконечно удаленной гиперплоскости.


Таким образом, результаты статьи \cite{A52} следуют из теоремы 1.


\medskip





{\bf 6.}


Рассмотрим нормализацию, дуальную центральной нормализации.


При такой нормализации все вторые нормали $K_x$


принадлежат фиксированной гиперплоскости $\alpha$.


Будем называть такую нормализацию {\it аффинной.}





\begin{thm}


Нормализация многообразия $X \subset \Bbb{P}^n$


является аффинной тогда и только тогда, когда $1$-формы


$\omega^0_p$ и $\omega^0_a$


в уравнениях $(17)$ обращаются в нуль\/${:}$


\begin{equation}
\omega_p^0 = 0, \quad \omega_a^0 = 0.
\end{equation}
Если нормализация многообразия $X \subset \Bbb{P}^n$ является


аффинной, то проективное пространство


$\Bbb{P}^n$ несет аффинную структуру, т.\,е.


$\Bbb{P}^n$ является аффинным пространством $\Bbb{A}^n$.


\end{thm}

\begin{pf}
Разместим точки $A_1, \dotsc, A_n$ подвижного репера на фиксированной


гиперплоскости $\alpha$. Поскольку при аффинной нормализации


$K_x \subset \alpha$, и, следовательно,


$dA_p \subset \alpha$, $p =l+1, \dotsc, n$, отсюда следует


$$


\omega_p^0 = 0.


$$
Более того, точки $A_a$, $a = 1, \dotsc, l$, первой нормали


$N_x$ также могут быть помещены в гиперплоскость $\alpha$.


Тогда $dA_a \subset \alpha$, откуда следует


$$


\omega_a^0 = 0.


$$

Обратно, если выполняется (32), то


$$
dA_p \subset \alpha, \quad d A_a \subset \alpha,
$$
где $\alpha = A_1 \wedge \dotsb \wedge A_n$.


Таким образом, плоскость


$\alpha$ является постоянной и нормализация многообразия


$X$ является аффинной, что и доказывает первую часть теоремы 2.





Для доказательства второй части теоремы 2 заметим, что


можно принять гиперплоскость~$\alpha$ за бесконечно удаленную


гиперплоскость $\Bbb{P}_\infty$ пространства $\Bbb{P}^n$.


Эта гиперплоскость определяет аффинную структуру в пространстве


$\Bbb{P}^n$.


Таким образом, пространство $\Bbb{P}^n$ становится аффинным


пространством $\Bbb{A}^n$.


\end{pf}





{\bf 7.} Предположим теперь, что нормализованное многообразие


$X \subset \Bbb{P}^n$ имеет плоскую нормальную связность


$\gamma_n$, т.\,е. $R^a_{bst} = 0$.


При этом из выражения (28) вытекают соотношения


\begin{equation}


\label{eq:43}


b^a_{pt} c^p_{bs} = b^a_{ps} c^p_{bt},
\end{equation}
которые только обозначениями отличаются от соотношений (3).


Используя матрицы


$$


B^a = (b^a_{pq}),\quad C_b = (c^p_{bq}),


$$
запишем соотношения (33) в виде


$$


(B^a C_b) = (B^a C_b)^T.


$$



Как доказано в (\cite{AG04}, гл.\,3 и 4), из этих соотношений


следует, что матрицы $B^a$ и $C_b$ могут быть одновременно приведены к


диагональному или блочно-диагональному виду.


Таким образом, доказана 



\begin{thm}
Фокусные гиперповерхности $F_x \subset N_x$ нормализованного


многообразия $X$ с плоской нормальной связностью


распадаются на плоские образующие различных размерностей.


\end{thm}

Это свойство многообразий $X$ с плоской нормальной связностью


$\gamma_n$ позволяет построить классификацию таких многообразий


таким же способом, как это было осуществлено для многообразий с


вырожденным гауссовым отображением в проективном пространстве.


Для многообразий в аффинном и евклидовом пространствах такая


классификация была описана в работах \cite{ACh75}--\cite{ACh01}.








\section{Нормализация многообразий в аффинном и евклидовом


пространствах}



{\bf 1.} Аффинное пространство $\Bbb{A}^n$ отличается от


проективного пространства


$\Bbb{P}^n$ тем, что в $\Bbb{A}^n$


фиксирована бесконечно удаленная гиперплоскость $\Bbb{P}_\infty$.


Если поместить вершины $A_i$, $i = 1, \dotsc, n$, подвижного


проективного репера в эту гиперплоскость, то уравнения


инфинитезимального перемещения подвижного репера принимают вид


(\cite{AG04}, уравнения (1.81))


\begin{equation}\label{eq:45}


\begin{cases}


dA_0 = \omega^0_0 A_0 + \omega^i_0 A_i, \\
dA_i \,= \phantom{\omega^0_0 A_0 +}\ \omega^j_i A_j,


\end{cases}


\qquad i, j = 1, \dotsc, n.


\end{equation}


При этом структурными уравнениями аффинного пространства являются


\begin{equation*}


%\label{eq:46}
d \omega^0_0 = 0, \ \  d \omega^i_0 = \omega_0^j \wedge
\omega_j^i, \ \  d \omega^i_j = \omega_j^k \wedge \omega_k^i.
\end{equation*}



Рассмотрим многообразие $X$ размерности $r$ в аффинном пространстве


$\Bbb{A}^n$.


 Касательное пространство $T_x(X)$ пересекает бесконечно удаленную


гиперплоскость $\Bbb{P}_\infty$ по подпространству $K_x$


размерности $r - 1$, $K_x = T_x \cap \Bbb{P}_\infty$.


 Таким образом, для нормализации многообразия


$X$ достаточно задать только семейство первых нормалей $N_x$.


 Если поместить точки $A_a, a = 1, \dotsc, l$,


подвижного репера в подпространство $N_x \cap \Bbb{P}_\infty$,


а точки $A_p$, $p = l+1, \dotsc, n$,


в подпространство $K_x$, то уравнения (34) примут вид


\begin{equation}\label{eq:47}


\begin{cases}
dA_0 = \omega^0_0 A_0\phantom{\omega_a^b A_b} + \omega^p_0 A_p, \\
dA_a =\phantom{\omega^0_0 A_0} \omega_a^b A_b +\omega^p_a A_p,\\
dA_p =\phantom{\omega^0_0 A_0} \omega_p^a A_a + \omega^q_p A_q
\end{cases}


\end{equation}
(см. уравнения (17)).





Как и в проективном пространстве, имеем уравнения (19),


где $b_{pq}^a$ --- второй фундаментальный тензор


многообразия $X$.


 Уравнения (21) также сохраняют свой вид,


однако вместо уравнений (20) получаем следующие уравнения:


\begin{equation}\label{eq:50}
\omega_p^0 = 0.
\end{equation}



Поскольку $l_{pq} = 0$, уравнения фокусной гиперповерхности


$F_x \subset N_x$ сохраняют свой вид (23).


Что же касается уравнения (24) фокусного гиперконуса $\Phi_x$,


то в соответствии с формулой (36)


это уравнение принимает вид


\begin{equation}\label{eq:52}
\det (\xi_a b^a_{pq}) = 0.
\end{equation}



Выражения (27) для компонент тензора кривизны аффинной связности


$\gamma_t$, индуцированной на нормализованном многообразии


$X\subset \Bbb{A}^n$, запишутся теперь в виде


\begin{equation}\label{eq:53}
R^p_{qst} = b^a_{qs} c_{at}^p - b^a_{qt} c_{as}^p.


\end{equation}


Что же касается выражений (28) для компонент тензора нормальной


кривизны многообразия $X$, то они сохраняются.





Рассмотрим тензор $R_{st}$, полученный из тензора кривизны $R^p_{qst}$


аффинной связности $\gamma_t$ сверткой по индексам $p$ и $q$.


Будем называть этот тензор


{\it тензором типа Риччи связности $\gamma_t$}.


Из выражения (38) следует


$$
R_{st} = b^a_{ps} c^p_{at} - b^a_{pt} c^p_{as}.
$$



Аналогичным образом определим


{\it тензор типа Риччи нормальной связности} $\gamma_n$.


Обозначим его $\widetilde{R}_{st}$.


Из выражений (28) получим


$$


\widetilde{R}_{st} = b^a_{pt} c^p_{as} - b^a_{ps} c^p_{at}.


$$
Сравнивая последние два уравнения, заключаем, что


$$


R_{st} = -\widetilde{R}_{st}.


$$



Таким образом, имеет место



\begin{thm}
На многообразии $X \subset \Bbb{A}^n$, снабженном аффинной


нормализацией, тензоры типа Риччи связностей


$\gamma_t$ и $\gamma_n$


отличаются лишь знаком.


\end{thm}



В случае нормализованной гиперповерхности $X$ в аффинном


пространстве $\Bbb{A}^n$ имеет место 





\begin{thm} Если аффинная связность $\gamma_t$,


индуцированная на нормализованной гиперповерхности $X \subset \Bbb{A}^n$,


является плоской, то и нормальная связность $\gamma_n$ является


плоской.


\end{thm}



\begin{pf} Действительно, в случае гиперповерхности $X$


используемые индексы пробегают следующие области значений:


$$


a, b = 1;\quad p, q, s, t = 2, \dotsc, n.


$$
Поэтому тензор кривизны нормальной связности $\gamma_n$ имеет


компоненты $R^1_{1st}$. Тогда из теоремы 4 следует


$$


R^1_{1st} = - R^p_{pst}.


$$
Однако если связность $\gamma_t$ является плоской, то $R^p_{qst} = 0$


и, следовательно, $R^p_{pst} = 0$.


 В результате имеем $R^1_{1st} = 0$, и поэтому связность


$\gamma_n$ также является плоской.


\end{pf}



Так же, как было в проективном пространстве, в аффинном пространстве


обращение в нуль тензора нормальной кривизны $R^a_{bst}$


эквивалентно полной интегрируемости


системы Пфаффа, определяющей распределение


$\Delta_y = y \wedge K_x$, где $y \in N_x$.


 Но в аффинном пространстве элементы $\Delta_y$


этого распределения параллельны подпространству $T_x (X)$.



Таким образом, имеет место 



\begin{thm}
 Нормализованное многообразие $X \subset \Bbb{A}^n$ имеет плоскую


нормальную связность $\gamma_n$ тогда и только тогда, когда это


многообразие допускает $l$-параметрическое


семейство параллельных многообразий $X (y)$, где $y \in N_x$.


\end{thm}






{\bf 2.} Другая связь теории многообразий с вырожденным


гауссовым отображением с теорией нормализованных многообразий была


установлена в \cite{Cha77}


(см. также \cite{Cha78}, теоремa 4;


\cite{Cha90}, теоремa~1, с.\,39).


Изложим эту теорему.



Предположим, что в каждой точке $x$ нормализованного многообразия


$X$ задано $s$-мерное направление $\nu^s (x)$ (т.\,е.


$s$-мерная плоскость, проходящая через $x$),


принадлежащее первой нормали $N_x (X)$.


 Это означает, что задано гладкое поле нормальных


$s$-мерных направлений $\nu^s (x)$ на $X$, где $s \leq l = n - r$.


 Это поле определяет нормальное подрасслоение $\nu^s (X)$, слоями


которого являются $s$-мерные центропроективные пространства.



 Плоскость $N^s (x)$ этого поля, соответствующая точке


$x \in X$, может быть задана точкой $x$ и точками


\begin{equation}\label{eq:55}
B_f = \xi_f^a A_a \in N_x,
\end{equation}
где $f, g, h = 1, \dotsc, s$.



Кроме этого, плоскость $N^s (x)$ должна быть инвариантна


по отношению к допустимым преобразованиям подвижного репера в $N_x (X)$.


Необходимые и достаточные условия ее инвариантности имеют вид


$$


dB_f = \theta_f^g B_g + \theta_f^0 A_0 \pmod{\omega^p},


$$
где $\theta_f^g$ и $\theta_f^0$ --- линейно независимые 1-формы.



Поле $\nu^s$ называется {\it параллельным} по отношению к


нормальной связности $\gamma_n$, если при любом инфинитезимальном


смещении произвольной точки $x \in X$ смещение


$s$-мерного направления $\nu^s (x)$ происходит в $(r + s)$-мерной


плоскости, натянутой на касательное подпространство $T_x


(X)$, $x \in X$, и направление $\nu^s (x)$.





Найдем аналитические условия параллельности поля $\nu^s$.


Всякое направление, принадлежащее $s$-мерному элементу


поля $\nu^s$, определяется точками $A_0 = x$ и


\begin{equation}\label{eq:57}
A = A_0 + \xi^f B_f,
\end{equation}
где $B_f$ имеют вид (39).



Дифференцируя соотношения (40) внешним образом и применяя (17),


получаем


\begin{equation*}


%\label{eq:58}
dA = (\omega_0^0 + \xi^f \xi_f^a \omega_a^0) A_0 + d \xi^f B_f +
(\omega^p + \xi^f \xi_f^a \omega_a^p) A_p + \xi^f (d\xi_f^a +
\xi_f^b \omega_b^a) A_a.
\end{equation*}
Поле $\nu^s$ параллельно в нормальной связности


$\gamma_n$ тогда и только тогда, когда


$$
(d\xi_f^a + \xi_f^b \omega_b^a) A_a = \theta_f^g B_g + \theta_f^0
A_0.
$$
Отсюда и из формулы (39) следует


$$


d\xi_f^a + \xi_f^b \omega_b^a = \theta^g_f \xi_g^a.


$$




\begin{thm}[{\series{m}\shape{n}\selectfont \cite{Cha77} или \cite{Cha90}, с.\,39--40}]


Поле $\nu^s$ \ $s$-мерных


нормальных направлений $\nu^s (x)$ на нормализованном многообразии


$X \subset \Bbb{P}^n$ является параллельным в нормальной связности


$\gamma_n$ тогда и только тогда, когда плоскости $N^s (x)$, $x \in X$,


образуют многообразие $V^{r+s}_{r}$ в $\Bbb{P}^n$


с вырожденным гауссовым


отображением ранга $r$ с $s$-мерными плоскостными образующими.


\end{thm}

Теорема 7 указывает метод построения многообразия $V^{r+s}_{r}$


с вырожденным гауссовым отображением, исходя из нормализованного


многообразия $X$ некоторого специального типа.





{\bf 3.} Рассмотрим, наконец, многообразие $X$ размерности $r$ в евклидовом


пространстве $\Bbb{E}^n$.


 На многообразии $X$ естественным образом определяются


вторая нормаль $K_x = T_x \cap \Bbb{P}_\infty$ и первая нормаль


$N_x$, ортогональная касательному пространству $T_x (X)$.



 В евклидовом пространстве $\Bbb{E}^n$ определено скалярное


произведение векторов.


 Это скалярное произведение индуцирует


скалярное произведение точек, принадлежащих


бесконечно удаленной гиперплоскости $\Bbb{P}_\infty$.


В рассматриваемом подвижном репере имеем


$A_a \in N_x \cap \Bbb{P}_\infty \overset{\mathrm{def}}{=}{} N'_x$;


$A_p \in T_x \cap \Bbb{P}_\infty = K_x$,


$a = 1, \dotsc, l$; $p = l + 1, \dotsc, n$;


и $T_x \perp N_x$.


Отсюда


\begin{equation}\label{eq:60}
(A_a, A_p) = 0,
\end{equation}
 где круглые скобки означают скалярное произведение точек бесконечно


удаленной гиперплоскости $\Bbb{P}_\infty$.


 Формулы


\begin{equation}\label{eq:61}


(A_a, A_b) = g_{ab}, \;\; (A_p, A_q) = g_{pq}


\end{equation}
определяют невырожденные симметрические тензоры $g_{ab}$ и $g_{pq}$.





Дифференцируя соотношения (41) и используя формулы


(35), (41) и (42), находим


$$


g_{ab}\, \omega_p^b + g_{pq}\, \omega_a^q = 0.


$$
Отсюда следует


$$


\omega_a^p = -g^{pq}\, g_{ab}\, \omega_q^b.


$$
Из уравнений (40) и (19), где $b^q_{pq}$  --- второй фундаментальный


тензор многообразия $X$, следует


\begin{equation}\label{eq:63}
\omega_a^p = -g^{pq}\, g_{ac} \, b_{qs}^c \omega^s.
\end{equation}


Сравнивая (43) и (21), получаем


\begin{equation}\label{eq:64}
c_{as}^p = -g^{pq} \,g_{ac} \,b_{qs}^c.
\end{equation}



В соответствии с (23) и (44) для фокусной гиперповерхности


$F_x$ многообразия $X \in \Bbb{E}^n$ имеем уравнение


$$


\det (y^0 \delta_q^p - y^a g^{ps} g_{ac} b_{sq}^c) = 0,


$$


которое эквивалентно уравнению


\begin{equation}\label{eq:65}
\det (y^0 g_{pq} - y_a b_{p q}^a) = 0,
\end{equation}
где $y_a = g_{ab} y^b$.





В рассматриваемом подвижном репере бесконечно удаленная


гиперплоскость $\Bbb{P}_\infty$ определяется уравнением $y^0 = 0$.


Поэтому в соответствии с (45) пересечение $F_x \cap \Bbb{P}_\infty$


фокусной гиперповерхности $F_x$ с бесконечно удаленной гиперплоскостью


$\Bbb{P}_\infty$ задается уравнением


\begin{equation}\label{eq:66}
\det (y_a b_{p q}^a) = 0.
\end{equation}


Но это уравнение лишь обозначениями отличается от уравнения


(37) фокусного гиперконуса $\Phi_x$ многообразия $X$.


Уравнения (37) и (46) совпадают, если $\xi_a = y_a = g_{ab} y^b$.





Последнее соотношение определяет поляритет в подпространстве $N'_x$


относительно мнимой квадрики $g_{ab} y^a y^b=0$, который точке


$y\in N'_x$ ставит в соответствие $(l-2)$-мерное подпространство


$\eta$, также принадлежащее подпространству $N'_x$.


 Подпространство $\eta$ вместе с подпространством $T_x$ определяет


гиперплоскость $\xi$, являющуюся линейной оболочкой подпространств


$T_x$ и~$\eta$.





В результате получена





\begin{thm}


Фокусный гиперконус $\Phi_x$ многообразия $X \subset \Bbb{E}^n$


образован гиперплоскостями $\xi$, которые являются линейными


оболочками подпространств $T_x$ и $\eta$, где $\eta\subset N'_x$ ---


$(l-2)$-мерные подпространства, полярные точкам $y\in F_x \cap N'_x$


относительно абсолюта пространства $N'_x$, определяемого уравнением


$g_{ab} y^a y^b=0$.


\end{thm}

Этот результат проясняет геометрический смысл фокусного гиперконуса


$\Phi_x$ для многообразия $X \subset \Bbb{E}^n$ и его связь с


фокусной гиперповерхностью $F_x$ многообразия $X$.





 Найдем еще тензор кривизны аффинной связности, индуцированной


на многообразии\linebreak $X \subset \Bbb{E}^n$.


 Для этого подставим в формулу (38)


значения $c_{aq}^p$, определяемое формулой (44).


В результате получаем


\begin{equation}\label{eq:67}
R^p_{qst} = g^{pu} g_{ac} (b^a_{qt} b^c_{us} - b^a_{qs} b_{ut}^c).
\end{equation}
Свертывая выражение (47) с тензором $g_{pv}$, имеем


\begin{equation}\label{eq:68}
R_{pqst} = g_{ac} (b^a_{ps} b^c_{qt} - b^a_{pt} b_{qs}^c),
\end{equation}


где $R_{pqst} = g_{pu} R^u_{qst}$.


Формулы (47) и (48) представляют собой обычные выражения тензора


кривизны аффинной связности $\gamma_t$, индуцированной на


многообразии $X \subset \Bbb{E}^n$.





Однако помимо тензора кривизны аффинной связности, индуцированной на


многообразии $X \subset \Bbb{E}^n$, мы


рассматривали также тензор нормальной кривизны $R^a_{bst}$, определенный


уравнениями (28).


Подставляя в формулу (28) значения $c_{aq}^p$ из (44), получим


\begin{equation*}


%\label{eq:69}
R^a_{bst} = g^{pq} g_{bc} (b^c_{qt} b_{ps}^a - b^c_{qs}
b_{pt}^a).
\end{equation*}





Как было отмечено выше, внешняя 2-форма


$$


\Omega_b^a = d \omega_b^a - \omega_b^c \wedge \omega_c^a =
\tfrac{1}{2}R^a_{bst} \omega^s \wedge \omega^t


$$
называется  в \cite{Ca01} {\it гауссовым кручением} многообразия


$X \subset \Bbb{E}^n$.




\begin{thebibliography}{99}





\bibitem{Fi50} 


Фиников С.П. {\em Теория конгруэнций}. --


М.--Л.: ГИТТЛ,  1950. -- 528 с.





\bibitem{AG04}


Akivis M.A., Goldberg V.V.


{\em Differential geometry of varieties with degenerate Gauss
maps}. -- New~York: Springer-Verlag, 2004. -- 276~p.



\bibitem{AG93} Akivis M.A., Goldberg V.V.


{\em Projective differential geometry of submanifolds}. --
Amsterdam: North-Holland, 1993. -- 374~p.





\bibitem{N76} Норден А.П.


{\it Пространства аффинной связности}. --


М.: Наука, 1976. -- 432~с.





\bibitem{KN}


Kobayashi S., Nomizu K.


{\it Foundations of differential geometry}. -- V.\,1. 


New York--London: Interscience


Publishers, John Wiley \& Sons, Inc., 1963. -- 340~p.





\bibitem{Ca01}


Cartan \'{E}. {\em  Riemannian geometry in an orthogonal
frame}. -- River Edge, NJ: World Sci. Publ. Co., Inc., 2001. --


277~p.





\bibitem{AG95}


Akivis M.A., Goldberg V.V.


{\it Normal connections of a submanifold in a projective
space} // Proc. of the Conference on Differential Geometry,
Hamiltonian Systems and Operator Theory (Univ. of West
Indies, Mona Campus, Jamaica, Feb.~7--11, 1994), 1995. --
P.\,137--158.



\bibitem{A52}


Атанасян Л.С.


{\it Оснащенные многообразия частного вида в многомерном


аффинном пространстве} // Тр. семин. по векторн. и


тензорн. анализу. -- М.: МГУ. -- 1952. -- Т.\,9. -- С.\,351--410.





\bibitem{ACh75}


Акивис М.А., Чакмазян А.В.


{\it Об oснащенных подмногообразиях


аффинного пространства, допускающих параллельное нормальное векторное


поле} // ДАН АрмССР. -- 1975. -- Т.\,60. -- \No\,3.


-- С.\,137--143.



\bibitem{ACh76}


Акивис М.А., Чакмазян А.В.


{\it О подмногообразиях евклидова пространства с плоской нормальной


связностью} //


ДАН АрмССР. -- 1976. -- Т.\,62. -- \No\,2. -- С.\,75--81.



\bibitem{ACh01}


Akivis M.A., Chakmazyan A.V.


{\it Dual-normalized submanifolds and hyperbands of curvature} //


Rend. Semin. matem. Messina. -- 2001--2002. -- Ser.~II. -- V.\,8. -- P.\,13--23.





\bibitem{Cha77}


Чакмазян А.В.


{\it Нормализованное по Нордену


подмногообразие в $P^n$


с параллельным нормальным подрасслоением} //


Матем. заметки. -- 1977. -- Т.\,22. -- \No\,5. -- С.\,649--662.



\bibitem{Cha78}


Чакмазян А.В.


{\it Связность в нормальных расслоениях нормализованного


подмногообразия $V^m$ в $P^n$} // Итоги науки и техн.


Пробл. геометрии. -- М.: ВИНИТИ, 1978. -- Т.\,10. -- С.\,55--74.





\bibitem{Cha90}


Чакмазян А.В.


{\it Нормальные связности в геометрии подмногообразий}. --


Ереван: Изд-во Армянск. гос. пед. инст., 1990. -- 116 с.





\end{thebibliography}





\vskip 1cm





{\small Примечание к пп.\,1, 5 и 6 литературы





1. Имеется немецкий перевод, сделанный Г.\,Болем: Akademie-Verlag,


1959. -- 506 p.





5. Имеется русский перевод, выполненный Л.В.Сабининым: М.: Наука, 1981. -- 344 с.





6. Эта книга является переводом на английский язык, выполненным


В.В.\,Гольдбергом, одноименной книги, изданной в 1960~г. на русском языке.


М.: Изд-во МГУ, 1960. -- 307 с.}





\vskip 1cm





\post{\it Иерусалимский технологический}{\it Поступила}


\post{\it колледж Махон-Лев $($Израиль$)$}{24.09.2003}


\smallskip


\post{\it Технологический институт}{}


\post{\it $($Ньюарк, штат Нью-Джерси, США$)$}{}


\smallskip


\post{\it Армянский государственный педагогический}{}


\post{\it институт $($Ереван, Армения$)$}{}





\label{end}


\end{document}


